\section{Limitations}

\label{sec:limitations}

The implementation and measurements target one platform:
\code{linux/arm64} on NVIDIA DGX Spark with \code{sm\_121}. They do not
establish portability or performance on x86-64, another GPU architecture, a
multi-GPU server, or a cloud accelerator. Two rounds on one host expose
run-to-run variation and reverse subject order, but they cannot characterize
population-level hardware variance or eliminate all persistent host-state and
thermal effects. Idle baselines and one-at-a-time CUDA isolation constrain
these effects without proving that they are absent.

The runtime supports only the pinned 1.7B, 12-Hz VoiceDesign checkpoint. It
does not implement voice cloning, reference-audio conditioning, speaker
enrollment, Base or CustomVoice checkpoints, or the 0.6B model. Explicit
language identifiers are limited to Chinese, English, Japanese, Korean,
German, French, Russian, Portuguese, Spanish, and Italian, plus automatic
selection. Unsupported languages are rejected rather than advertised from
anecdotal output.

The benchmark measures serving behavior, not speech naturalness,
intelligibility, instruction adherence, accent, or listener preference. A
faster result would not by itself demonstrate better audio. Conversely, the
same released checkpoint does not guarantee bit-identical waveforms across two
different implementations unless a separate parity test establishes that
property.

The native and stock HTTP paths do not expose identical termination semantics.
Native reports natural EOS explicitly; stock raw PCM does not. The evaluation
therefore preserves Stock EOS as unknown even when transport completion and a
conservative length gate succeed. Finally, the current native path does not use
CUDA Graph capture or replay, so this paper makes no claim about their effect.
